\section{Version 1: Simple Edge Deletion} \label{sec:kps_v1}

The first version of the generic algorithm relies on a straightforward approach
to break the cycles. To obtain the set of vertex-disjoint paths $P$ from the
cycle cover $\mathcal{C}$, the algorithm simply deletes the minimum-weight edge
from each cycle.

In the absolute worst-case scenario—where the cycle cover $\mathcal{C}$ consists
entirely of 2-cycles and both edges in every 2-cycle have identical weights
(i.e., $b = c$)—the weight of the resulting paths $w(P)$ can be as small as
$\frac{1}{2}W$. However, in most instances, the presence of larger cycles and
weight imbalances ensures a significantly better result.

To analyze this formally, let $pW$ denote the total weight of the set consisting
of the lightest edge from each $3^+$-cycle in $\mathcal{C}$. Because we remove
exactly the lighter edge from each 2-cycle (totaling $cW$) and the lightest edge
from each $3^+$-cycle (totaling $pW$), the total weight of the remaining paths
is precisely:
\begin{equation}
w(P) = (1 - c - p)W
\end{equation}

By definition, the lightest edge in any $3^+$-cycle can contain at most
one-third of the total weight of that cycle. Since the total weight of all
$3^+$-cycles in $\mathcal{C}$ is exactly $(1 - b - c)W$, we can bound $p$ by the
inequality:
\begin{equation} \label{eq:p_bound}
p \le \frac{1}{3}(1 - b - c)
\end{equation}

By substituting the upper bound from Inequality \ref{eq:p_bound} into the
equation for $w(P)$, we can establish the guaranteed approximation ratio for
this version.

\begin{theorem} \label{lemma:version1}
Version 1 of the generic algorithm returns a $\left(\frac{2}{3} + \frac{1}{3}(b
- 2c)\right)$-approximation to the directed Max TSP.
\end{theorem}

\begin{proof}
Substituting $p \le \frac{1}{3}(1 - b - c)$ into the weight function gives:
\begin{equation}
w(P) \ge \left(1 - c - \frac{1}{3}(1 - b - c)\right)W = \left(\frac{2}{3} + \frac{1}{3}(b - 2c)\right)W
\end{equation}
Since the weight of the optimal tour $w(T_{opt})$ is bounded from above by the
weight of the optimal cycle cover $W$, the resulting set of paths $P$ guarantees
the stated approximation ratio.
\end{proof}

\subsection{Computational Complexity}
Let $n$ denote the number of vertices of $G$. Computing the maximum-weight cycle
cover $\mathcal{C}$ via the Hungarian algorithm takes $O(n^3)$ time
\cite{kuhn1955hungarian, edmondskarp1972, tomizawa1971techniques}. Since the cycles of $\mathcal{C}$ partition the vertex
set, deleting the minimum-weight edge from each cycle -- and thereby obtaining
the set of paths $P$ -- takes only $O(n)$ time overall, as does patching the
paths of $P$ into a single tour. The cycle cover computation therefore
dominates, and Version 1 runs in $O(n^3)$ time.
